\begin{lemma}
For every $q\in\mathbb N$, the hypergraph $A_q$ of
\noderef{OrderThreeAffineMultihypergraph} belongs to $H(3,3,3q+1)$.
In its line graph an index is adjacent to $\ast$ if and only if it differs
from $\ast$, and $\deg(\ast)=9q$.
The edges indexed by $(m,a,i)$ and $(n,b,j)$ are disjoint if and only if
$m=n$ and $a\ne b$.
\end{lemma}
\begin{proof}
Use the explicit edge sets of \noderef{OrderThreeAffineMultihypergraph}.
The parametrizations $y\mapsto(0,y)$ and $x\mapsto(x,mx+a)$ are injective,
so every edge has three vertices. This gives \noderef{UniformHypergraph}.
An intersection is a subset of either edge, hence has size at most three,
which gives \noderef{TSimpleHypergraph}, including for different copies of
the same line.

Fix a vertex $(x,y)$. Among nonvertical indices incident with it the map
$(m,i)\mapsto(m,y-mx,i)$ is a bijection from
$(\mathbb Z/3\mathbb Z)\times\{0,\ldots,q-1\}$: the equation
$y=mx+a$ uniquely determines $a=y-mx$, and forgetting $a$ is its inverse.
There are therefore $3q$ such indices. The vertical index is incident
exactly when $x=0$, so the degree is $3q+1$ or $3q$.
By \noderef{HypergraphDegree} and \noderef{MaxDegreeAtMost} this proves
the required degree bound; \noderef{HypergraphClass} combines the three properties.

Every nonvertical edge contains $(0,a)$, which belongs to the vertical
edge. By \noderef{LineGraphOfHypergraph}, each nonvertical index is adjacent
to $\ast$, and $\ast$ is not adjacent to itself. Its neighbors are in
bijection with triples $(m,a,i)$, of which there are $3\cdot3\cdot q=9q$.

If $m=n$ and $a\ne b$, a common point would imply $mx+a=mx+b$,
contradicting cancellation. If $m=n$ and $a=b$, the lines coincide and
contain $(0,a)$, so are not disjoint. Finally, if $m\ne n$, the nonzero
residue $m-n$ is invertible modulo three. Set $x=(m-n)^{-1}(b-a)$
and $y=mx+a$. Then $(m-n)x=b-a$ implies $y=nx+b$, so $(x,y)$
is a common point. These cases prove the asserted equivalence.
\end{proof}
